\documentclass[tikz,border=1pt]{standalone}
\usepackage{amsmath,amssymb,bm}
\usepackage{tikz}
\usetikzlibrary{decorations.markings,hobby,topaths}

\newcommand{\MakeStrandCoords}[5]{\pgfmathtruncatemacro{\Ns}{#2}
  \def\Phase{0}\ifx#5B\def\Phase{0.5}\fi
  \foreach \i in {0,...,\Ns}{
    \pgfmathsetmacro{\s}{\i/\Ns}
    \pgfmathsetmacro{\y}{#3*sin(360*(#4*\s + \Phase))}
    \path[postaction={decorate, decoration={markings,
      mark=at position \s with {\coordinate (#5\i) at (0,\y);}}}] #1;
  }}
\newcommand{\drawSeg}[4]{\ifnum#4=1
    \draw[preaction={draw=\Core, line width=\Mask}, fat strand=#3]
      (#1#2) .. (#1\the\numexpr#2+1\relax);
  \else
    \draw[fat strand=#3]
      (#1#2) .. (#1\the\numexpr#2+1\relax);
  \fi}
\newcommand{\Amp}{0.13}
\newcommand{\Turns}{4}
\newcommand{\Samples}{400}
\newcommand{\Wth}{2.5pt}
\newcommand{\Core}{black}
\newcommand{\Mask}{5.0pt}
\newcommand{\ClrA}{black!60!black}
\newcommand{\ClrB}{red!70!black}
\newcommand{\KPATH}{([closed] P1)..(P2)..(P3)..(P4)..(P5)..(P6)..(P7)..(P8)..(P9)}
\newcommand{\Phase}{0}

\usepackage{tikz}
\usetikzlibrary{decorations.markings,hobby}

\usetikzlibrary{
    spath3,
    intersections,
    arrows,
    knots,
    calc,
    hobby,
    decorations.pathreplacing,
    shapes.geometric,
}

% ---------------- Global styles (safe for 'knots') ----------------
\tikzset{
    knot diagram/every strand/.append style={
        line cap=round,
        line join=round,
        ultra thick,
        black
    },
    every knot/.style={line cap=round,line join=round,very thick},
    strand/.style={line cap=round,line join=round,line width=3pt,draw=black},
    over/.style={preaction={draw=white,line width=6.5pt}},
% DO NOT set a global "every path" here; it breaks internal clip paths.
}
% ------- TikZ Guide Lines -------
\newcommand{\SSTGuidesPoints}[2]{% #1=basename (e.g. P), #2=last index
    \foreach \i in {1,...,#2}{
      \fill[blue] (#1\i) circle (1.2pt);
      \node[blue,font=\scriptsize,above] at (#1\i) {\i};
    }
    \draw[gray!40, dashed]
    \foreach \i [remember=\i as \lasti (initially 1)] in {2,...,#2,1} { (#1\lasti)--(#1\i) };
}
\providecommand{\swirlarrow}{\rightsquigarrow}
\begin{document}
\begin{tikzpicture}[use Hobby shortcut, line cap=round, line join=round, scale=1.05]
\def0.13{0.13}
\def4{4}        
\def400{400}
\def2.5pt{2.5pt}
\defblack{black}     
\def5.0pt{5.0pt}
\defblack!60!black{black!60!black}
\defred!70!black{red!70!black}

\tikzset{
  fat strand/.style={
    draw=#1, line width=2.5pt,
    double=black, double distance=1.2*2.5pt
  }
}


\coordinate (P1) at (-2.0,-2.0);
\coordinate (P2) at (-2.0, 2.0);
\coordinate (P3) at ( 1.0,-0.5);
\coordinate (P4) at (-1.0,-0.5);
\coordinate (P5) at ( 2.0, 2.0);
\coordinate (P6) at ( 2.0,-2.0);
\coordinate (P7) at (-1.0, 0.5);
\coordinate (P8) at ( 1.0, 0.5);
\coordinate (P9) at (-2.0,-2.0); 
\def([closed] P1)..(P2)..(P3)..(P4)..(P5)..(P6)..(P7)..(P8)..(P9){([closed] P1)..(P2)..(P3)..(P4)..(P5)..(P6)..(P7)..(P8)..(P9)}


\newcommand{\MakeStrandCoords}[5]{
  \pgfmathtruncatemacro{\Ns}{#2}
  \def0{0}\ifx#5B\def0{0.5}\fi
  \foreach \i in {0,...,\Ns}{
    \pgfmathsetmacro{\s}{\i/\Ns}
    \pgfmathsetmacro{\y}{#3*sin(360*(#4*\s + 0))}
    \path[postaction={decorate, decoration={markings,
      mark=at position \s with {\coordinate (#5\i) at (0,\y);}}}] #1;
  }
}
\MakeStrandCoords{([closed] P1)..(P2)..(P3)..(P4)..(P5)..(P6)..(P7)..(P8)..(P9)}{400}{0.13}{4}{A}
\MakeStrandCoords{([closed] P1)..(P2)..(P3)..(P4)..(P5)..(P6)..(P7)..(P8)..(P9)}{400}{0.13}{4}{B}


\newcommand{\drawSeg}[4]{
  \ifnum#4=1
    \draw[preaction={draw=black, line width=5.0pt}, fat strand=#3]
      (#1#2) .. (#1\the\numexpr#2+1\relax);
  \else
    \draw[fat strand=#3]
      (#1#2) .. (#1\the\numexpr#2+1\relax);
  \fi
}


\pgfmathtruncatemacro{\Last}{400-1}
\foreach \k in {0,...,\Last}{
  \pgfmathsetmacro{\s}{\k/400}
  
  \pgfmathtruncatemacro{\blk}{floor(2*4*\s + 0.5)}
  \ifodd\blk
    
    \drawSeg{A}{\k}{black!60!black}{0}
    \drawSeg{B}{\k}{red!70!black}{1}
  \else
    
    \drawSeg{B}{\k}{red!70!black}{0}
    \drawSeg{A}{\k}{black!60!black}{1}
  \fi
}
\end{tikzpicture}
\end{document}